\section{RAE for multi-agent systems}
\subsection{OMPAS: Extending RAE for fleet management}

\begin{frame}{The Operational Model Planning and Acting System (OMPAS), an extension to RAE}
    \begin{columns}[t]
        \begin{column}{0.5\textwidth}
    New task in separated thread
            \tikzstyle{thread} = [draw, rounded corners, rectangle, fill=black!50]
    \begin{tikzpicture}[scale=0.5, every node/.style={scale=0.5}]
  
        %RAE
        \node[draw = black, very thick,
            minimum width = 14em,
            minimum height = 5em,
            rounded corners,
            text = black,
            text centered, text depth = 4 cm,
            align = center,
            ] (R) at (0,0) {\textbf{OMPAS}};
        %Operational Models
        % \node[draw = black, ellipse, very thick,
        %     minimum width = 7em,
        %     minimum height = 2em,
        %     above = 1em of R,
        %     text = black,
        %     text centered, align=center
        % ] (OM) {\textit{Operational Models} %($A_\Delta$)
        % };
  
  
        %state
        \node[draw = black, very thick,
        text centered,
        minimum width = 13em,
        minimum height = 2em,
        rounded corners,
        text = black,
        fill=black!10,
        align = center] (S) at ($(R.south) + (0,0.6)$) {State};
  
        %main
        \node[draw = black, very thick,
        minimum height = 8em,
        minimum width = 13em,
        rounded corners,
        text = black,
        text centered, text depth = 7em,
        fill=black!10,
        align=center] (main) at ($(R.north) + (0,-2)$) {Main};
  
        %t1
        \node[draw = black, thick,
        minimum height = 6em,
        minimum width = 7em,
        rounded corners,
        text = black,
        text centered, text depth = 5em,
        fill=black!20,
        align=center](t1) at ($(main.south) + (-0.8, 1.3)$) {\emph{transport($o_1$,$l_7$)}};
  
  
        %first task
        \node[draw, rounded corners, rectangle, fill=black!30] (st11) at ($(t1.south) + (0,1.3)$) {\emph{drive($t_1$,$l_8$)}};
        \node[draw, rounded corners, rectangle, fill=black!30] (st12) at ($(st11.south) + (0,-.5)$){\emph{pack($o_1$)}};
        
        \node[draw, thick, rounded corners, rectangle,minimum width = 4em, fill=black!20, align= center] (t3) at ($(t1.east) + (0.9, 0.3)$){\textit{monitor}\\
        \textit{-fuel}($t_1$)};
        \node[draw, thick, rounded corners, rectangle,minimum width = 4em, fill=black!20, align= center] (t3) at ($(t1.east) + (0.9, -0.5)$){\dots};
  
  
  
        %platform
        \node[draw = black, very thick,
        text centered,
        minimum width = 5em,
        minimum height = 2em,
        %below = 2em of R,
        rounded corners,
        text = black,
        align = center,] (Pl) at ($(R.west) +(-8em, -2em)$) {Platform};
  
        \node[text centered,
        minimum width = 5em,
        minimum height = 2em,
        below = 1.5em of Pl,
        rounded corners,
        text = black,
        align = center,] (E) {\textbf{\large\textit{Environment}}};
  
        \node[very thick,
            text centered,
            rounded corners,
            text = black] (U) at ($(R.west) +(-8em, +4em)$) {\Large \textbf{User}};
        
        % \node[draw = black, very thick,
        %     text centered,
        %     right = 7em of R,
        %     minimum width = 5em,
        %     minimum height = 4em,
        %     rounded corners,
        %     text = black] (Pl) {Platform};
        
  
        \path[->]
  
            %Link between operational models library
            %(OM) edge (R)
  
            (U.20) edge[bend left] node [midway, above, align=center] {\footnotesize \emph{Tasks}} (main.+140)
            (R.+160) edge [bend left] node [midway, above] {\footnotesize \emph{Reports}} (U.-20)
        
            %Link between RAE and the platform
            (main.-160) edge[bend left] node [pos = 0.6, above = 0.2em] {\footnotesize \emph{Commands}} (Pl.east)
            
            %between platform and state
            (Pl.-20) edge [bend right] node [pos = 0.7, right, above, align = center] {\footnotesize \emph{Updates}} (S.west)
            
            %between platform and main
            (Pl.20) edge [bend left] node [pos = 0.3, right, above= 0.2em, align = center] {\footnotesize \emph{Events}} (main.-170)
  
            (Pl.-110) edge (E.110)
            (E.70) edge (Pl.-70)
            ;
        
        \end{tikzpicture}
        \end{column}
        \begin{column}{0.5\textwidth}
        Acting language with support for
        \begin{itemize}
            \item Generic programming constructs
            \item Concurrency
            \item Error and interruption
        \end{itemize}
        \end{column}
    \end{columns}
\end{frame}

\subsection{Acting language}

\begin{frame}[fragile]{An acting language for concurrency and automated analysis of skills}
    
    Language requirements :
    \begin{itemize}
        \item Generic programming constructs,
        \item Acting features:
        \begin{itemize}
            \item control of robotic platform
            \item explicit choice points
        \end{itemize}
        \item Support for Concurrency
        \begin{itemize}
            \item Interruption handling
            \item Resource modelling
        \end{itemize}
        \item Identified and simple semantic $\rightarrow$ automated analysis of skills to guide choice
    \end{itemize}

\end{frame}

\begin{frame}[fragile]{SOMPAS: A Lisp variant based on Scheme}
    \begin{columns}[t]

        \begin{column}{0.4\textwidth}
            Scheme advantages:
            \begin{itemize}
                \item Few primitives
                \item Functional
                \item Simple to extend
            \end{itemize}
            Definitions:
            \begin{itemize}
                \item LValue: Atom, List
                \item Atom: Symbol, Boolean, Number, Procedure
            \end{itemize}
        \end{column}
        \begin{column}{0.6\textwidth}
            Example:
                        \lstset{columns=fullflexible}
                        \begin{lstlisting}
(begin
    (define loc-o (read-state loc o))
    (define loc-t (read-state loc t))
    (if (!= loc-o loc-t)
        (exec drive t loc-o))
    (exec load t o) 
    (exec drive t l)
    (exec unload t o))             
                        \end{lstlisting}
                    \end{column}
     \end{columns}

\end{frame}

\begin{frame}{Scheme extension for concurrency and interruption}

    Definition of a new Atom: the \emph{handle}
    
    ~~

    Functions:
    \begin{itemize}
        \item \textit{(async e)}: creates a new thread to evaluate e, returns a handle
        \item \textit{(await h)}: awaits the result of the concurrent evaluation.
        \item \textit{(interrupt h)}: interrupts the concurrent evaluation.
    \end{itemize}
\end{frame}


% \begin{frame}[fragile]{Example of program with concurrency and interruptions.}
    
% \begin{columns}
%     \begin{column}{0.4\textwidth}
%         \lstset{basicstyle=\footnotesize, columns=fullflexible}
% \begin{lstlisting}[language=lisp]
% (begin
%  (define h
%   (async 
%     (begin
%      (uninterruptible (begin
%        (exec pick ?r ?o)
%        (exec move ?r ?l)
%        (exec drop ?r ?o)))
%      (exec inspect ?r ?o))))
%  (race 
%    (await h)
%    (begin
%     (sleep 10)
%     (interrupt h))))
% \end{lstlisting}

%     \end{column}
%     \begin{column}{0.6\textwidth}
%         \newcommand{\Cross}{$\mathbin{\tikz [x=1.4em,y=1.4em,line width=.2em, red] \draw (0,0) -- (1,1) (0,1) -- (1,0);}$}
%     \begin{tikzpicture}
%         \node[draw,
%         rounded corners,
%         ultra thick,
%         minimum height = 3em,
%             minimum width = 7em,] (pick) at (0,0) {pick r1 o3};
    
%         \node[draw, rounded corners,
%         ultra thick,
%             minimum height = 3em,
%             minimum width = 7em,
%         ] (move) at ($1.5*(pick.south)-0.5*(pick.north)$) {move r1 l6};
%         %\action{move, move r1 l6, 10em, pick}
    
%         \node[draw, rounded corners,
%         ultra thick,
%             minimum height = 3em,
%             minimum width = 7em,
%         ] (drop) at ($1.5*(move.south) - 0.5*(move.north)$) {drop r1 o3};
%         \node[draw, rounded corners,
%         ultra thick,
%             minimum height = 3em,
%             minimum width = 7em,
%         ] (inspect) at ($1.5*(drop.south) - 0.5*(drop.north)$) {inspect r1 o3};
%         \node at (inspect) {\Cross};
    
%             \node[] (origin) at ($(pick.north west) + (-1em,+0.5em)$) {~~};
        
%             \node[] (zero) at ($(origin.west) + (-1em,0)$) {0};
%             \draw[-, ultra thick] (origin.west) -- (origin.east);
%             \node[] (ten) at ($(zero) + (0,-6em)$) {10};
%             \draw[-, ultra thick] ($(origin.west) + (0,-6em)$) -- ($(origin.east) + (0,-6em)$);
    
    
%             \node[] (interruption) at ($(ten.east) +(11em,0em)$){\color{red} \textit{Int.}};%{\color{red} \textit{interruption}};
%             \draw[-, ultra thick, dashed, color= red] ($(ten.east)$) -- ( interruption);
%             \node[] (end) at ($1.5*(inspect.south west) - 0.5*(inspect.south west) + (-1em,-1em)$) {};
%             \node[] (time) at ($1.5*(end.south east) - 0.5*(end.south east) + (2em,0)$)  {time (seconds)};
    
    
%         \path[->] (origin.center) edge[ultra thick] (end);
%     \end{tikzpicture}
%     \end{column}
% \end{columns}    
% \end{frame}

\begin{frame}{Acting primitives in SOMPAS}
    \begin{itemize}
    \pause
        \item \textit{(exec a $p_1...p_n$)} executes and monitors a task or command.
    \pause
        \begin{itemize}
            \item command : resorts to the platform
            \item task: selects a method and execute its operational model.
        \end{itemize}
    \pause
        \item \textit{(read-state sf $p_1...p_n$)} returns the value of a state-variable
        \item \textit{(arbitrary set $\lambda$)} returns an arbitrary element from a $set$. $\lambda$ can be used to select among the $set$.
    \end{itemize}
\end{frame}

\begin{frame}[t]{SOMPAS: Resources}
    Definition of a resource:
    \begin{itemize}
        \item Object with initial capacity $C(0) = C_{init} $
        \item acquire($r$,$c$) at time $t$: $c \leq C(t)$ 
    \end{itemize}

    ~~

    Kinds of resources : 
    \begin{itemize}
        \item \emph{unary}: $C_{init} = 1, c = 1$;
        \item \emph{divisible}: $C_{init} \in  \mathbb{R}_+^*, c \in ]0, C(t)]$
    \end{itemize}

    ~~

    Functions:
    \begin{itemize}
        \item \textit{(new-resource r C)} : create a new resource $r$ with capacity $C$
        \item \emph{(acquire r c)} : acquisition request for $r$ with amount $c$, returns a resource handle
        \item \emph{(release h)} : release the resource using the resource handle
    \end{itemize}

\end{frame}

\begin{frame}[fragile]{Multi-agent model}
    \begin{columns}
        \begin{column}{0.5\textwidth}
            \centering
            \includegraphics[width = 0.3\textwidth]{images/godot/package.png}
            
            \Large $\times n$
        \end{column}
        \begin{column}{0.5\textwidth}
            \centering
            \includegraphics[width = 0.3\textwidth]{images/godot/robot_texture.png}
            
            \LARGE $\times m$
        \end{column}
    \end{columns}

~~

        Modelled with a unique high-level task:
        \setlength{\leftmargini}{0pt}

        \tiny
    \begin{lstlisting}
(do
    (mapf new-resource (instance robot))
    (mapf new-resource (instance machine))
    (define h1 (async (t_process_packages)))
    (define h2 (async (t_check_rob_bat)))
    (await h1))
    \end{lstlisting}    
\end{frame}

\begin{frame}{Running example for multi-agent}
    \centering
    \LARGE
    \href{https://youtu.be/KqBnKKIjug4}{Video example}
\end{frame}

% \begin{frame}{More advanced features and programming constructs}
%     More complex behavior definition
%     \begin{itemize}
%         \item \textit{(wait-for dyn)} waits until the expression \textit{dyn} becomes true 
%         \item \textit{(run-monitoring e dyn)} evaluates \textit{e} while \textit{dyn} is true, interrupts \textit{e} otherwise.
%         \item ...
%     \end{itemize}
% \end{frame}

% \begin{frame}[t,fragile]{Defining an acting domain with SOMPAS}
%     \setlength{\leftmargini}{0pt}
%     \begin{columns}
%         \begin{column}{0.5\textwidth}
%             \begin{itemize}
%                 \footnotesize
%                 \item Task (label, typed parameters):
%                 \tiny
%                 \begin{lstlisting}
% (def-task t_process_package
%     (:params (?p package)))
%                 \end{lstlisting}
%                 \footnotesize
%                 \pause
%                 \item Method (label, typed parameters, pre-conditions, body):
%                 \tiny
%                 \begin{lstlisting}
% (def-method m_process_to_do_r
%   (:task t_process_package)
%   (:params (?p package))
%   (:pre-conditions
%     (!= (package.processes_list ?p)
%         nil)))
%   (:body
%     (do
%      (define ?m ...)
%      (t_process_on_machine ?p ?m)
%      (t_process_package ?p)))   
%                 \end{lstlisting}
%             \end{itemize}
%         \end{column}
%     \pause
%         \begin{column}{0.5\textwidth}
%             \begin{itemize}
%     \pause
%             \footnotesize
%             \item State-function (label, typed parameters \& result):
%             \tiny
%             \begin{lstlisting}
% (def-state-function at
%     (:params (?r robot))
%     (:result location))
%             \end{lstlisting}
%     \pause
%             \footnotesize
%             \item Command (label, typed parameters):
%             \tiny
%             \begin{lstlisting}
% (def-command pick (:params (?r robot)))
%                 \end{lstlisting}
%             \end{itemize}

%     \pause
%             ~
%         \normalsize
%         Others:
        
%         \textit{lambda, type, constant}
%         \end{column}

%     \end{columns}
% \end{frame}